\documentclass[../../../include/open-logic-section]{subfiles}

\begin{document}

\olfileid[id]{sth}{ord-arithmetic}{mult}

\olsection{Perkalian Ordinal}

Sekarang kita beralih ke perkalian ordinal, yang kita dekati dengan cara yang
sangat mirip dengan penjumlahan ordinal. Andaikan kita hendak mengalikan
$\alpha$ dengan~$\beta$. Untuk melakukannya, bayangkan suatu kisi persegi
panjang dengan lebar $\alpha$ dan tinggi $\beta$; hasil kali $\alpha$ dan
$\beta$ diperoleh dengan bergerak sepanjang setiap baris, lalu bergerak
melalui baris berikutnya\ldots{} sampai seluruh kisi telah dilalui. Dengan
kata lain, hasil kali $\alpha$ dan $\beta$ diperoleh dengan mengganti
\emph{setiap} anggota dalam $\beta$ dengan suatu salinan~$\alpha$.

Untuk memformalkan gagasan ini, kita cukup menggunakan pengurutan
leksikografis terbalik pada hasil kali Kartesius $\alpha$ dan~$\beta$:

\begin{defn}
Untuk sebarang ordinal $\alpha, \beta$, hasil kalinya ialah
$\alpha \ordtimes \beta = \ordtype{\alpha \times \beta, \rlexless}$.
\end{defn}
\noindent
Sekali lagi, kita harus memastikan bahwa definisi ini terbentuk dengan baik:

\begin{lem}\ollabel{ordtimeslessiswo}
$\tuple{\alpha \times \beta, \rlexless}$ merupakan pengurutan baik bagi
sebarang ordinal $\alpha$ dan~$\beta$.
\end{lem}

\begin{proof}
Tepat seperti pada \olref[add]{ordsumlessiswo}.
\end{proof}
\noindent
Tidak sulit pula membuktikan bahwa perkalian berperilaku sebagai berikut:

\begin{lem}\ollabel{ordtimesrecursion}
Untuk sebarang ordinal $\alpha, \beta$:
\begin{align*}
	\alpha \ordtimes 0 &= 0\\
	\alpha \ordtimes (\beta \ordplus 1) &=
		(\alpha \ordtimes \beta) \ordplus \alpha\\
	\alpha  \ordtimes \beta &=
		\supstrict_{\delta < \beta}(\alpha \ordtimes \delta) &&
		\text{jika $\beta$ merupakan ordinal limit}.
\end{align*}
\end{lem}

\begin{proof}
Diserahkan sebagai latihan.
\end{proof}

Sesungguhnya, sebagaimana dalam kasus penjumlahan, kita dapat mendefinisikan
perkalian ordinal melalui persamaan-persamaan rekursi ini alih-alih memberikan
definisi langsung. Seperti halnya penjumlahan, beberapa perilakunya juga
terasa akrab:

\begin{lem}\ollabel{ordinalmultiplicationisnice}
Jika $\alpha, \beta, \gamma$ ordinal, maka:
\begin{enumerate}
	\item\ollabel{ordtimes1} jika $\alpha \neq 0$ dan $\beta < \gamma$,
	maka $\alpha \ordtimes \beta < \alpha \ordtimes \gamma$;
	\item\ollabel{ordtimes2} jika $\alpha \neq 0$ dan $\alpha \ordtimes
	\beta = \alpha\ordtimes\gamma$, maka $\beta = \gamma$;
	\item\ollabel{ordtimes3} $\alpha \ordtimes (\beta \ordtimes
	\gamma) = (\alpha \ordtimes \beta) \ordtimes \gamma$;
	\item\ollabel{ordtimes4} jika $\alpha \leq \beta$, maka $\alpha
	\ordtimes \gamma \leq \beta \ordtimes \gamma$;
	\item\ollabel{ordtimes5} $\alpha \ordtimes (\beta \ordplus
	\gamma) = (\alpha \ordtimes \beta )\ordplus (\alpha\ordtimes
	\gamma)$.
\end{enumerate}
\end{lem}

\begin{proof}
Diserahkan sebagai latihan.
\end{proof}

Anda dapat membuktikan (atau mencari) hasil-hasil lain sesuka hati. Akan
tetapi, mengingat \olref[ord-arithmetic][add]{ordsumnotcommute}, hasil berikut
seharusnya tidak mengejutkan:

\begin{prop}
Perkalian ordinal tidak komutatif: $2 \ordtimes \omega =
\omega < \omega \ordtimes 2$.
\end{prop}

\begin{proof}
$2 \ordtimes \omega = \supstrict_{n < \omega}(2\ordtimes n) = \omega \in
\supstrict_{n < \omega}(\omega \ordplus n) = \omega \ordplus \omega =
\omega \ordtimes 2$.
\end{proof}
\noindent
Sekali lagi, alasan intuitifnya cukup lugas. Untuk menghitung $2 \ordtimes
\omega$, kita mengganti setiap bilangan asli dengan dua entitas. Tipe urutan
yang sama akan diperoleh jika kita cukup menyisipkan semua bilangan
``setengah'' ke dalam bilangan asli, yakni jika kita mempertimbangkan
pengurutan alami pada $\Setabs{\nicefrac{n}{2}}{n \in \omega}$. Dengan
ungkapan ini, jelas bahwa tipe urutannya sama dengan tipe urutan $\omega$
sendiri. Namun, untuk menghitung $\omega \ordtimes 2$, kita menempatkan dua
salinan $\omega$, yang satu sesudah yang lain.

\begin{prob}
Buktikan
\olref[sth][ord-arithmetic][mult]{ordtimeslessiswo},
\olref[sth][ord-arithmetic][mult]{ordtimesrecursion},
dan
\olref[sth][ord-arithmetic][mult]{ordinalmultiplicationisnice}.
\end{prob}

\end{document}
